\needspace{10\baselineskip}
\SourceLeaf{316}
\begin{center}
\pdfbookmark[1]{14. Correction of the elements from four oppositions}{zusatz-014}
[14.]
\end{center}
\begin{center}
\{On correcting the elements of a planet\\from four observed oppositions.\}
\end{center}
If the elements of a planet are already known so accurately that the adopted values differ very little from the true ones, and observations of four oppositions are available, the elements can be corrected more readily and simply than if the four given positions of the planet had been chosen arbitrarily. The theory of this investigation has already been set out in the paper »Disquisitio de elementis ellipticis Palladis«. Here follow a few practical remarks applicable to the calculation of oppositions and the like in general.

First the epoch must be corrected. Calculate the geocentric longitude and latitude of the planet from any observation, preferably one belonging to one of the middle oppositions. From these, using the radius vector obtained by the known formulas from the unchanged epoch, calculate the heliocentric longitude and latitude, and from these the longitude in the orbit. After applying the perturbations to this longitude (if formulas for them are already available), derive the mean anomaly from it. Comparison with the mean anomaly obtained from the unchanged epoch gives the change in the epoch. If no observations are given and only the calculated longitudes and latitudes at the time of the four oppositions are available, one of these longitudes must be chosen. It is better, however, to make the observations oneself than to rely on other people's calculations.

Now calculate the four oppositions provisionally using the epoch corrected in this way and keeping the other elements unchanged. Earlier calculations can be used here with sufficient accuracy. The most convenient procedure for calculating an opposition is as follows.

If observations from several places are available, select the observations from the place that supplied the most. From the stated times of transit across the meridian, determine the times of culmination of the planet at two-day intervals throughout the period covered by all the observations. In a preliminary calculation for a day in the middle of the observation period, determine the distance of the planet from the Earth, and from this the time it takes light to travel from the planet to the Earth. This time must be subtracted from the times of culmination, thereby allowing for aberration. Now derive the mean anomaly from the corrected epoch by subtracting the longitude of perihelion at the time of the epoch. From this, using the daily tropical motion, determine the mean anomalies for the chosen days, and from these the right ascensions and declinations. To allow for nutation, use the true longitude of the node, obtained by adding the nutation in longitude (\textsc{Bohnenberger}'s Astron. p. 672) to the mean longitude given by the elements. Apply the parallax according to the method explained in the Theoria motus Corpor. Coel. art. 59 or 70. It is easy to see, moreover, why the calculation is carried out at two-day intervals even when no observations were made on some of those days. Differences provide a very convenient check on the calculation, so that no error can escape notice, and precise interpolation then gives the calculations for \textls[150]{all} days.

Compare the right ascensions and declinations thus obtained with those observed and determine from them the mean deviations in right ascension and in declination. By adding these deviations to the calculated values, derive the geocentric longitude for the day of opposition, which is very soon found by \SourceLeaf{317}comparing the observations with the calculated solar longitudes. Determine the same longitude at a time differing by an arbitrary amount (say a few hours) from the time previously assumed, and then use interpolation to find the instant at which the longitudes of the Sun and the planet differ by exactly \(180^{0}\). In the calculation for the Sun, moreover, find the \textls[150]{true} longitude at the time of culmination corrected for the aberration of the planet, that is, add the Sun's aberration to the longitude obtained from the tables, which already includes it.

If only a few observations are available (say up to six), one may also calculate the right ascensions and declinations for the observation days alone, if that seems preferable.

In the preliminary calculation of the oppositions one may, if necessary, neglect aberration, nutation and parallax. If the oppositions are to be calculated only once, these must be taken into account.

From the longitudes and latitudes found at the time of the four oppositions, derive the corrections to the individual elements by the method given in the paper de Elementis Palladis.

\par\medskip\begin{center}\rule{.22\textwidth}{.3pt}\end{center}
